\section{System Architecture}
The Vertex-9 core architecture utilizes a decoupled design topology, isolating the host interface domain from the high-frequency DSP processing pipelines. All system control, DMA configuration, and signal parameters are governed by the control and status register file.

\begin{tcolorbox}[
    colback=lightbg,
    colframe=bordergray,
    title=\textbf{Figure 1: Internal FPGA Architecture Block Diagram},
    coltitle=primary,
    fonttitle=\sffamily\bfseries\small,
    arc=3mm,
    left=10pt,
    right=10pt,
    top=10pt,
    bottom=10pt
]
\centering
\begin{tikzpicture}[
    box/.style={draw, rectangle, rounded corners=2pt, minimum width=2.2cm, minimum height=1.0cm, align=center, thick, fill=white, draw=primary, font=\sffamily\scriptsize},
    bus/.style={draw, ->, >=stealth, ultra thick, secondary},
    signal/.style={draw, ->, >=stealth, thick, accent, dashed},
    every node/.style={font=\sffamily\scriptsize},
    scale=0.92,
    every node/.append style={scale=0.92}
]
  % Nodes
  \node[box] (pcie) at (0,3) {\textbf{ExpressLink Gen3}\\\textbf{Controller}\\Hard IP Block};
  \node[box] (dma) at (3.5,3) {\textbf{Scatter-Gather}\\\textbf{DMA Engine}\\Vertex Core};
  \node[box] (regs) at (7,3) {\textbf{CSR Interface}\\\textbf{\& Registers}\\SynapseBus Slave};
  \node[box] (dsp) at (7,0.5) {\textbf{DSP Pipeline}\\\textbf{Processor}\\Parallel IQ Core};
  
  \node[box, fill=accent!5, draw=accent] (ddr) at (11,3) {\textbf{DDR4 Controller}\\\textbf{\& Arbiter}\\SynapseBus Mem};
  \node[box, fill=accent!5, draw=accent] (phy) at (11,0.5) {\textbf{High-Speed SerDes}\\\textbf{ADC/DAC Interface}\\Vertex Phy};
  
  % Arrows
  \draw[bus] (pcie) -- (dma) node[midway, above] {V-Stream};
  \draw[bus] (dma) -- (regs) node[midway, above] {SynapseBus-L};
  \draw[bus] (dma) -- (ddr) node[midway, above, sloped] {V-Bus Master};
  \draw[bus] (phy) -- (dsp) node[midway, above] {IQ Signals};
  \draw[bus] (dsp) -- (ddr) node[midway, left] {V-Stream};
  
  % Control signals
  \draw[signal] (regs) -- (dsp) node[midway, left] {Ctrl};
  \draw[signal] (regs) -- (dma) node[midway, left] {Config};
  
  % External labels
  \draw[<->, thick] (-1.8, 3) -- (pcie) node[midway, above] {PCIe Lanes};
  \draw[<->, thick] (ddr) -- (13.2, 3) node[midway, above] {DDR4 Memory};
  \draw[<->, thick] (phy) -- (13.2, 0.5) node[midway, above] {RF Frontend};

\end{tikzpicture}
\end{tcolorbox}

\subsection{Core Subsystem Descriptions}
\begin{itemize}
    \item \textbf{ExpressLink Gen3 Controller:} Fictional hard macro handling physical, link, and transaction layer protocols for host interfacing. Bridges standard system PCIe commands to internal protocols.
    \item \textbf{Scatter-Gather DMA Engine:} Dedicated controller designed to support ring-buffered, non-contiguous host memory access. Coordinates dual-channel memory-mapped transfers to and from DDR4 RAM.
    \item \textbf{Control \& Status Registers (CSR):} Interconnect endpoints accessed via the SynapseBus-Lite register protocol. Manages hardware parameter configuration and reports processing health.
    \item \textbf{DSP Pipeline Processor:} Fixed-point signal processor consisting of parallel multiply-accumulate (MAC) arrays and circular pipeline registers.
\end{itemize}

% Section 3: Register Map
